\documentclass[12pt]{article}

\usepackage[a4paper,margin=1in]{geometry}
\usepackage{setspace}
\usepackage{amsmath}
\usepackage{enumitem}

\setstretch{1.2}

\begin{document}

\section{Method}

\subsection{Model setup}

\subsubsection{Genomic architecture and selection regime from Hayward (2022)}
\begin{enumerate}
    \item Highly polygenic QTL for the trait under selection
    \item Population size decreased (supplement validation figure 1)
    \item Fitness function is wide, though scaled by the reduced population size
    \item Shift of the optimum is a constant and a simulation parameter
\end{enumerate}

\subsubsection{Derived phenotype calculations based on Moore et al.\ (1987)}
\begin{enumerate}
    \item Direct implementation of the Moore model works in SLiM (supplement validation figure 2)
    \item Redefine $\psi$ because of not translational invariance
    \item Three models of phenotype calculations with examples
\end{enumerate}

\subsection{Analytical solution}

\subsubsection{Analytically derived genotypic and phenotypic distributions}
\begin{enumerate}
    \item Copy the average
    \item Copy the extreme
    \item Copy the best (TODO)
\end{enumerate}

\subsubsection{Invasion analysis}
\begin{enumerate}
    \item Copy the average
    \item Copy the extreme
    \item Copy the best (TODO)
\end{enumerate}

\section{Results}

\subsection{Single change in environment}

\subsubsection{Figure 1: Genotypic and phenotypic distribution under stabilizing selection (generation 9999)}
\begin{enumerate}
    \item Show that the genotypic distribution is normal
    \item Phenotypic distribution differs depending on $\psi$, group size, and model, without changes in the genotypic distribution
\end{enumerate}

\subsubsection{Figure 2 / Supplemental Figure 1: Distribution of the first generation under directional selection (generation 10000)}
\begin{enumerate}
    \item Show the phenotypic response to the changed optimum
    \item Response depends on $\psi$, group size, and model
    \item Reason for supplement: main point is the shift in mean phenotype, shown in adaptation trajectories and first-generation phenotype jump heat map
    \item Alternative plot: scatter plot of mean phenotype change as a function of group size to show diminishing returns
\end{enumerate}

\subsubsection{Figure 3: Adaptation trajectories}
\begin{enumerate}
    \item Show slower movement of mean phenotype toward the new optimum with increasing sociality
    \item Including genotype trajectories can show slowdown of genotypic evolution due to masking/shielding
    % \item Alternatively, plot differences between mean genotype and mean phenotype over time to show genotypic catch-up
\end{enumerate}

\subsubsection{Figure 4: Heat map of first phenotypic jump for copy-the-best}
\begin{enumerate}
    \item Show dependence of initial response on $\psi$ and group size with diminishing returns
    \item Analytical explanation of plateau near 25 for group size=50:
    \begin{itemize}
        \item Optimum shift of 120 is far in the tail of the fitness function
        \item Copy-the-best effectively becomes copy-the-max
        \item Plateau value approximated by
        \[
            \psi \times (\text{genetic standard deviation}) \times E[\max(a_1, a_2, \ldots, a_n)]
        \]
    \end{itemize}
\end{enumerate}

\subsubsection{Figure 5: Difference in decay constants}
\begin{enumerate}
    \item Compare $\psi = 0$ and $\psi = 0.9$ as a function of group size
    \item Show slowdown of adaptation affected by group size and $\psi$
\end{enumerate}

\subsection{Fluctuating environment}

\subsubsection{Figure 1: Distribution right before an environmental shift}
\begin{enumerate}
    % \item Show genotypic distribution centered between two optima ()
    \item Distribution is not centered on either individual optimum
\end{enumerate}

\subsubsection{Figure 2: Phenotypic trajectory over a full cycle}
\begin{enumerate}
    \item Include normalized fitness trajectory
    \item Intermediate $\psi$ (e.g.\ $\psi=0.5$) may yield higher fitness than nonsocial case at the end of the shift
\end{enumerate}

\subsubsection{Figure 3: Summary statistics}
\begin{enumerate}
    \item Some stats shown as a function of group size and $\psi$
    \item could be describing the exponential fit to the trajectory or the distribution
\end{enumerate}

\subsection{Supplementary Figures}

\subsubsection{Validation figures}
\begin{enumerate}
    \item Distance to the new optimum normalized by initial genetic variance at the shift
    \item Compare $\psi=0$, $N=1000$ simulation to Hayward (2022)
    \item Direct implementation of Moore et al.\ model under truncated selection
    \item Show SLiM simulations reproduce Moore model single-generation predictions
\end{enumerate}

\end{document}
